\section{Introduction}

Large language models pass classical logic benchmarks at impressive rates,
but it is not obvious whether they do so by carrying internal representations
of logical structure or by matching surface statistics. The distinction is
methodologically important: only the former predicts generalisation to novel
formal languages.

We adopt the working hypothesis from the mechanistic interpretability
literature \cite{wang2022interpret,nanda2023progress}: that components of a
trained transformer can be localised and that their roles can be tested by
selective perturbation. We specialise the question to logic:

\begin{quote}
\emph{Does the model carry a linearly decodable representation of validity
that survives a change of surface form to an invented symbolic language?}
\end{quote}

We answer affirmatively for GPT-2 with two complementary diagnostics. First,
per-layer logistic-regression probes trained on either English or a
controlled conlang reach high accuracy; the linear separability is highest
in the middle blocks (layer~$\gptTwoProbeEnPeakLayer$ for English,
layer~$\gptTwoProbeClPeakLayer$ for conlang). Second, cross-surface
probes---train on one surface, test on the other---transfer at a rate well
above chance, confirming that the validity direction is not entirely
language-specific. Finally, we exhaustively ablate every attention head and
report layer-wise sensitivity profiles.

\paragraph{Contributions.}
\begin{enumerate}
\item A reproducible parallel English/conlang benchmark of propositional
and syllogistic logic, balanced across templates and labels.
\item Per-layer probe accuracies on both surfaces, with cross-surface
transfer scores.
\item Exhaustive head-ablation diagrams contrasting impact on logic
accuracy against impact on language modelling.
\end{enumerate}
